% Draft 2: interpretation-led prose, annotator terminology and cited judge method.
% One anonymous human reference sheet; authored edit types are not reference labels.
\section{Results on New Data}
\label{sec:new-results}

\subsection{Predictive Performance}

The three variants (Natural, Ordinary and Homoscedastic) were evaluated on the additional public datasets and synthetic regimes introduced in Section~\ref{sec:new-data}.
The learner and training budget remain unchanged from the replication.
Stage selection follows the controlled replication's validation-NLL rule rather than the released script's squared-error rule.
Homoscedastic uses the revised scale calculation described in Section~\ref{sec:new-data}.
Parkinsons and grouped Airfoil additionally test predictions for unseen subjects or conditions.

\begin{table*}[t]
\centering
\small
\setlength{\tabcolsep}{4pt}
\begin{tabular}{@{}lrrrrrrrrr@{}}
\toprule
& \multicolumn{3}{c}{Natural} & \multicolumn{3}{c}{Ordinary} & \multicolumn{3}{c}{Homoscedastic} \\
Setting & NLL & RMSE & $C_{95}$ & NLL & RMSE & $C_{95}$ & NLL & RMSE & $C_{95}$ \\
\midrule
Airfoil, rows & 2.2470 & 2.592 & 0.913 & 2.9465 & 5.414 & 0.956 & \textbf{2.1567} & 2.073 & 0.957 \\
Airfoil, conditions & 2.5453 & 3.441 & 0.938 & 3.1586 & 5.928 & 0.917 & \textbf{2.4559} & 2.719 & 0.956 \\
Forest Fires & \textbf{1.7785} & 1.410 & 0.935 & 1.7826 & 1.404 & 0.933 & 1.7787 & 1.415 & 0.935 \\
Parkinsons, subjects & \textbf{3.6301} & 8.572 & 0.919 & 3.6486 & 8.358 & 0.924 & 3.6662 & 9.000 & 0.932 \\
Real Estate & \textbf{3.6441} & 8.059 & 0.959 & 4.0367 & 13.506 & 0.971 & 3.7928 & 7.955 & 0.959 \\
Superconductivity, full & \textbf{3.5158} & 11.558 & 0.945 & 4.7108 & 33.858 & 0.990 & 3.7974 & 10.762 & 0.942 \\
\bottomrule
\end{tabular}
\caption{Additional public-data results: unweighted means of fold or repetition metrics under the partitions in Table~\ref{tab:new-data-overview}.
Lower NLL and RMSE are better; $C_{95}$ is the fraction of test outcomes inside the nominal 95\% prediction interval, with 0.95 the nominal target rather than a score to maximise.
Bold identifies the lowest mean NLL within each setting, not a significance claim.
RMSE retains each dataset's target units; Forest Fires uses the transformed target.
NLL magnitudes are not comparable across target scales.}
\label{tab:new-public-results}
\end{table*}

\textbf{Public-data findings.}
Natural has lower mean test NLL than Ordinary in all 6 settings, agreeing with Yacht's direction but reversing the Kin8nm ordering in both the published and local replication results (Tables~\ref{tab:new-public-results} and~\ref{tab:replication-results}).
Natural achieves the lowest NLL on Parkinsons, Real Estate and full Superconductivity, while Homoscedastic performs best on both Airfoil evaluations.
Forest Fires shows little numerical separation between the variants.
NLL and RMSE give different rankings in four of the six settings, agreeing only on the Airfoil evaluations: accurate point predictions do not necessarily identify the best predictive distributions.

\textbf{Interpreting coverage.}
On Superconductivity, Ordinary's highest coverage comes with mean interval widths of 115.79 against 34.55 for Natural, and with worse NLL and point predictions, so coverage alone does not indicate better probabilistic performance (Table~\ref{tab:new-public-results}; Appendix~\ref{app:new-results-details}).

\textbf{Generalisation and training budget.}
All three variants have higher mean NLL when Airfoil conditions are held out together than under row holdouts (Table~\ref{tab:new-public-results}).
This tests prediction for unseen condition combinations, although changed training proportions mean that the difference cannot be attributed to grouping alone.
For Superconductivity, repeated predictor combinations can cross partitions, so the results concern new rows rather than necessarily new materials.
Moreover, 12 of its 15 fits select the maximum stage count: these are fixed-budget results, not evidence that further training cannot help.
Appendix~\ref{app:new-results-details} records aggregation details and cap counts.

\textbf{Constructed-data findings.}
Natural has the lowest NLL in the input-dependent-noise regimes A and C, while Homoscedastic leads in the constant-noise control B (Table~\ref{tab:synthetic-results}).
The paired replicate comparisons support these orderings and extend Natural's Yacht advantage over Ordinary across all three regimes (Appendix~\ref{app:new-results-details}).
Homoscedastic has the highest NLL in the local Yacht and Kin8nm replication runs, but the lowest in B.
Its revised scale calculation means that this reversal cannot be attributed to the dataset alone.

\begin{table}[t]
\centering
\small
\setlength{\tabcolsep}{3.5pt}
\begin{tabular}{@{}llrrr@{}}
\toprule
Regime & Variant & NLL & Scale MAE & Oracle gap \\
\midrule
A & Natural & \textbf{0.8249} & 0.063 & 0.077 \\
& Ordinary & 0.8563 & 0.089 & 0.109 \\
& Homoscedastic & 1.0379 & 0.259 & 0.290 \\
\midrule
B & Natural & 0.9456 & 0.043 & 0.032 \\
& Ordinary & 0.9916 & 0.089 & 0.078 \\
& Homoscedastic & \textbf{0.9364} & 0.012 & 0.023 \\
\midrule
C & Natural & \textbf{1.4876} & 0.107 & 0.047 \\
& Ordinary & 1.5164 & 0.136 & 0.076 \\
& Homoscedastic & 1.7220 & 0.516 & 0.281 \\
\bottomrule
\end{tabular}
\caption{Means over 10 fresh synthetic replicates per regime.
Scale MAE compares predicted standard deviation with the true conditional noise scale.
Oracle gap is fitted NLL minus NLL under the true generating distribution on the same test outcomes; it is a reference difference, not a guaranteed nonnegative bound.
Bold marks the lowest mean NLL per regime.}
\label{tab:synthetic-results}
\end{table}

The scale errors help explain this contrast: Natural recovers the varying noise more accurately in A and C, whereas a single scale is more accurate when the noise is constant in B (Table~\ref{tab:synthetic-results}).
Natural and Homoscedastic nevertheless have nearly identical RMSE and overall 95\% coverage in A and C, for example 0.6820 against 0.6821 and 0.935 against 0.934 in A (Appendix~\ref{app:new-results-details}).
Because the true noise scales are known, the synthetic data show that models with similar point predictions and coverage can spread uncertainty differently across inputs.
This supports input-dependent scale estimation in the varying-noise settings tested here.
These findings concern simulated outcomes that are Normal given the inputs.
Regimes A and B differ in average noise variance as well as noise structure, so their comparison does not isolate noise structure alone.

\subsection{LLM-as-Judge Evaluation}

A separate LLM-as-judge experiment evaluates explanation clarity, interval interpretation and evidence support, complementing numerical assessment of predictive distributions.
\citet{zheng2023judging} study language models as evaluators.
The judge assesses written explanations, not regression variants.

\textbf{Evaluation design.}
MiniCPM4.1-8B \cite{openbmb2025minicpm41}, using Q8\_0 quantisation and served locally through Ollama~0.9.6 \cite{ollama2025}, evaluated 15 held-out Natural predictions, with 3 cases each from synthetic A, B and C, grouped Airfoil and Real Estate.
Each case supplied a prediction packet, one intended-faithful description and one controlled edit.
Across the 15 cases, five edits targeted interval misinterpretation, five unsupported causes and five obscuring jargon.
Packets omitted observed outcomes; synthetic packets additionally supplied the known generating law.
Implementation checks and a three-case smoke test preceded evaluation; the same prompt was then held fixed, with temperature 0 and seed 0, to measure the judge's behaviour.
All 60 requests produced valid structured outputs: 30 pointwise assessments on three binary criteria and 30 pairwise comparisons covering both answer orders.
Clarity assesses understandable wording, interval consistency assesses numbers, units and statistical interpretation, and evidence support assesses whether substantive claims follow from the packet.
A human annotator assessed all 30 explanations using the supplied evidence and rubric (Appendix~\ref{app:new-results-details}).
The presented items contained neither judge verdicts nor the authored edit categories.
These annotations are the agreement reference; the edit categories are not reference labels.

\begin{table}[t]
\centering
\small
\setlength{\tabcolsep}{4pt}
\begin{tabular}{@{}lrrr@{}}
\toprule
Criterion & Agreement & $\kappa$ & Baseline \\
\midrule
Clarity & 90.0\% & 0.000 & 90.0\% \\
Interval consistency & 86.7\% & 0.636 & 83.3\% \\
Evidence support & 93.3\% & 0.867 & 50.0\% \\
\bottomrule
\end{tabular}
\caption{Pointwise judge agreement with one human annotator over 30 explanations per criterion.
Baseline always predicts that criterion's majority human label; $\kappa$ is Cohen's chance-adjusted agreement \cite{cohen1960agreement}.}
\label{tab:judge-human-results}
\end{table}

\textbf{Agreement and failure modes.}
Evidence support agrees most strongly with the human reference, while clarity's 90\% agreement equals its majority baseline because the judge passes every explanation (Table~\ref{tab:judge-human-results}).
All disagreements on interval consistency and evidence support are explanations that the annotator accepted and the judge failed (Appendix~\ref{app:new-results-details}).
On three jargon-heavy explanations, the annotator identifies a clarity problem while accepting the interval statements; the judge assigns the opposite labels.
This suggests that the judge does not consistently separate readability from statistical correctness.
Annotator clarity labels also vary across the five items sharing the same final jargon sentence.

\textbf{Answer-order sensitivity.}
Only 4 of 15 cases retain the same underlying pairwise verdict after the explanations are swapped, and the judge favours the second position overall (Appendix~\ref{app:new-results-details}), consistent with the type of position bias documented by \citet{zheng2023judging}.
In all five interval-misinterpretation cases, it prefers the explanation that confuses a prediction interval with a confidence interval and claims a guarantee for the next outcome, whenever that explanation appears second.
Yet the same judge rejects all five texts on interval consistency when assessing them individually.
Pointwise agreement therefore does not establish reliable pairwise explanation selection.

Agreement was measured against one human annotator and describes this evaluation set; variation between annotators was not assessed.
Because the examples share five explanation templates, the results describe the tested examples rather than broad performance across explanation styles.
